\clearpage
\fancyhead[L]{\small JUNO.}
\hypertarget{native-261-lower}{}
\begin{center}
JUNO.

\rule{0.75\linewidth}{0.4pt}

Monatliche Correspondenz. Band X. S. 463–471. November 1804.

\rule{0.75\linewidth}{0.4pt}
\end{center}

\noindent—\quad—\quad—\par
Dr. \textsc{Gauss} bat uns, die Position eines Sterns zu bestimmen, der am 28. Sept. sehr nahe bei \emph{Juno} stand und mit dem er sie verglichen hatte. Aus dem Mittelwert der Beobachtungen an fünf Tagen erhielten wir die folgende Position dieses Sterns, die auch anderen Beobachtern nützlich sein kann.

\begin{center}
\begingroup\fontsize{9}{11}\selectfont\setlength{\tabcolsep}{3pt}\renewcommand{\arraystretch}{1.1}
\begin{tabular}{r|r|r|r|r}
Größe & \shortstack{Mittlere Rektaszension\\1804} & \shortstack{Jährliche\\Änderung} & \shortstack{Mittlere südliche\\Deklination\\1804} & \shortstack{Jährliche\\Änderung}\\ \hline
8 & $358^\circ\ 19'\ 34''42$ & $+46''1$ & $5^\circ\ 13'\ 57''27$ & $-20''0$ \\
\end{tabular}\endgroup
\end{center}


Unermüdlich setzt Dr. \textsc{Gauss} die beschwerliche und mühevolle Arbeit fort, die Elemente der Bahn dieses neuen Planeten zu verbessern. „Seit meinem letzten Brief“, (schreibt dieser große Rechner) „habe ich unsere \emph{Juno} weiterhin fleißig beobachtet. Für die Elemente verwende ich meine eigenen Beobachtungen aber nicht mehr, da Ihre Beobachtungen, die Sie mir freundlicherweise mitgeteilt haben, schon einen beträchtlichen Bogen zu umfassen beginnen. Inzwischen habe ich bei einer erneuten Verbesserung der Juno-Bahn vergeblich versucht, die mittlere Bewegung von Ceres und Pallas auch für Juno hypothetisch zugrunde zu legen. Darauf habe ich die Elemente ein zweites Mal und unabhängig von Hypothesen bestimmt. Dabei ergaben sich folgende Werte.“

\begin{center}
\emph{II\textsuperscript{te} Elemente der Juno.}
\end{center}

\begin{center}
\begingroup\fontsize{9}{11}\selectfont\setlength{\tabcolsep}{3pt}\renewcommand{\arraystretch}{1.1}
\begin{tabular}{lr}
Epoche 1804, 30. Sept., um $0^{\mathrm u}$ mittlerer Zeit  &  \\
\quad in Seeberg . . . . . . . . . & $21^\circ\ 17'\ 47''$ \\
Tägliche mittlere Bewegung . . . . & $836''89$ \\
Aphel . . . . . . . . . & $231^\circ\ 38'\ 1''$ \\
Exzentrizität . . . . . . . . . & $0.254964$ \\
Logarithmus der Halbachse . . . . . & $0.4182255$ \\
Knoten . . . . . . . . . . . & $170^\circ\ 46'\ 41''$ \\
Neigung der Bahn . . . . . . . & $12^\circ\ 19'\ 43''$ \\
\end{tabular}\endgroup
\end{center}


„Wenn sich in meine Rechnungen keine Fehler eingeschlichen haben, kann es schon als ziemlich sicher gelten, dass Junos mittlere Bewegung der von Ceres und Pallas \emph{nicht gleich} ist. Sie ist \emph{beträchtlich größer}, und folglich sind die Umlaufzeit und die mittlere Entfernung von der Sonne kleiner. Die weiteren Beobachtungen werden uns sehr bald Gewissheit darüber geben.“

Anfangs schien es, als würde Dr. \textsc{Gauss} auch für \emph{Juno} dieselbe Umlaufzeit wie für Ceres und Pallas finden. Das hätte der Hypothese von Dr. \textsc{Olbers} über die Entstehung dieses kleinen Plane%
\SourceLeaf{262}%
ten sehr widersprochen. Setzt man nämlich voraus, dass all diese kleinen Planeten vielleicht nur Trümmer und Bruchstücke eines zerstörten größeren Planeten sind, so lässt sich zwar zeigen, dass die Umlaufzeiten dieser kleinen Trümmer nicht \emph{sehr verschieden} sein können. Dennoch war es sehr unwahrscheinlich, dass alle Stücke des zersprengten Planeten genau dieselbe Geschwindigkeit erhalten haben sollten. Nun findet unser vortrefflicher \textsc{Gauss}, wie wir oben gesehen haben, dass die mittlere Bewegung von \emph{Juno} viel rascher ist als die von Ceres und Pallas.

Dr. \textsc{Olbers} äußert sich dazu folgendermaßen. „Die gesamte Lage der Juno-Bahn ist mit meiner Hypothese vereinbar (die ich übrigens weiterhin nur als Hypothese vorbringen will). Ihre Knoten mit der Ceres-Bahn liegen jetzt etwa 24 Grad vom Knoten der Pallas-Bahn entfernt. Bei den schon so unterschiedlichen Neigungen dieser Bahnen müssen sich die Knoten durch die Anziehungskraft des Jupiter aber ungleichförmig verschieben. Jetzt liegt die Juno-Bahn beim absteigenden Knoten auf der Ceres-Bahn, wo die Pallas-Bahn dieser so nahe kommt, weit innerhalb der Ceres-Bahn. Die Aphelien aller dieser Bahnen bewegen sich ganz anders als die Knoten. Die Lagen der Apsidenlinien gegenüber den Knoten verändern sich also ständig. Da diese Bahnen fast gleich große Achsen, aber sehr unterschiedliche Exzentrizitäten haben, folgt daraus, dass sie sich zu bestimmten Zeiten tatsächlich schneiden werden und sich auch in der Vergangenheit tatsächlich geschnitten haben. Nehme ich zum Beispiel die von \textsc{Oriani} bestimmte jährliche Verschiebung der Aphelien an, für Pallas $106''1$ und für Ceres $120''9$, und setze ich voraus, dass die Knoten gegenüber den Fixsternen ruhen und die Neigungen unveränderlich sind, so folgt, dass sich die Bahnen von Ceres und Pallas beim absteigenden Knoten von Pallas auf der Ceres-Bahn vor 7463 Jahren tatsächlich geschnitten haben und nach 282 Jahren wieder schneiden werden. Beim aufsteigenden Knoten wird ein solcher Schnitt in 925 Jahren eintreten. Wie jetzt die Pallas-Bahn an beiden Knoten innerhalb der Ceres-Bahn liegt, wird nach 1000 Jahren die Ceres-Bahn innerhalb der Pallas-Bahn liegen. Diese Überlegungen können jedoch erst dann zu einem entscheidenden Ergebnis führen, wenn die Störungen aller drei Bahnen vollständig ausgearbeitet sind.“

\Needspace{24\baselineskip}
Bis das geschehen kann, verbessert Dr. \textsc{Gauss} weiterhin die elliptischen Elemente der Bahn und vergleicht sie mit den besten Meridianbeobachtungen. Es folgt ein solcher Vergleich aller dieser Beobachtungen mit seinen zweiten Elementen der \emph{Juno}.

\begin{center}
\begingroup\fontsize{8}{10}\selectfont\setlength{\tabcolsep}{3pt}\renewcommand{\arraystretch}{1.1}
\begin{tabular}{r|r|r|r|r|l}
 & \multicolumn{2}{c|}{Berechnete} & \multicolumn{2}{c|}{Differenz} & \\ 1804 & \shortstack{Rektaszension} & \shortstack{Südliche\\Deklination} & \shortstack{in der\\Rektaszension} & \shortstack{in der\\Deklination} & Beobachter \\ \hline
13. Sept. & $0^\circ\ 44'\ 56''0$ & $1^\circ\ 52'\ 37''6$ & $-0''4$ & $+0''6$ & von Zach \\
14 & $0\ 35\ 42.9$ & $2\ 5\ 40.3$ & $+1.9$ & $+4.8$ & — \\
15 & $0\ 26\ 19.8$ & $2\ 18\ 48.6$ & $-0.2$ & $+10.1$ & — \\
17 & $0\ 7\ 8.0$ & $2\ 45\ 18.8$ & $-0.1$ & $-5.7$ & — \\
18 & $359\ 57\ 20.9$ & $2\ 58\ 38.7$ & $+0.4$ & $+7.8$ & — \\
20 & $359\ 37\ 39.0$ & $3\ 25\ 24.7$ & $.\ .\ .\ .$ & $+21.2$ & — \\
23 & $359\ 7\ 17.4$ & $4\ 5\ 36.5$ & $+2.5$ & $+2.0$ & — \\
23 & $359\ 7\ 3.3$ & $4\ 5\ 55.0$ & $+3.3$ & $+8.0$ & Dr. Burckhardt \\
25 & $358\ 46\ 45.8$ & $4\ 32\ 41.4$ & $-2.2$ & $+13.2$ & Dr. Maskelyne \\
27 & $358\ 26\ 56.9$ & $4\ 58\ 43.7$ & $-3.7$ & $-12.8$ & von Zach \\
28 & $358\ 16\ 58.7$ & $5\ 11\ 49.3$ & $-7.9$ & $+5.5$ & — \\
29 & $358\ 6\ 48.0$ & $5\ 25\ 11.7$ & $-4.3$ & $+15.0$ & Dr. Maskelyne \\
30 & $357\ 57\ 18.2$ & $5\ 37\ 50.4$ & $-8.5$ & $+6.7$ & von Zach \\
2. Okt. & $357\ 38\ 6.4$ & $6\ 2\ 59.7$ & $-9.6$ & $+0.2$ & — \\
4 & $357\ 19\ 34.8$ & $6\ 27\ 40.2$ & $-4.7$ & $+1.5$ & — \\
5 & $357\ 10\ 37.3$ & $6\ 39\ 43.7$ & $-4.2$ & $+3.8$ & — \\
6 & $357\ 1\ 52.8$ & $6\ 51\ 35.6$ & $-3.5$ & $+3.3$ & — \\
\end{tabular}\endgroup
\end{center}

\SourceLeaf{263}
\indent Nach denselben Elementen hat Dr. \textsc{Gauss} ferner für Mitternacht in Seeberg die folgende Ephemeride des künftigen Laufs der \emph{Juno} berechnet. [Vom 15. Okt. bis zum 2. Dez. enthält sie Rektaszension und Deklination.]

Etwa am 18. Dezember kommen sich Juno und Ceres geozentrisch nahe.

\noindent—\quad—\quad—\par

\hfill[\textsc{von Zach.}]

